\documentclass[aspectratio=169,11pt]{beamer}

% ------------------------------------------------------------
% Presentación: IA como aliada del emprendedor
% XXVII Semana de la Facultad de Economía, UAEMéx
% ------------------------------------------------------------
\usepackage[utf8]{inputenc}
\usepackage[T1]{fontenc}
\usepackage[spanish,es-nodecimaldot]{babel}
\usepackage{lmodern}
\usepackage{microtype}
\usepackage{amsmath,amssymb}
\usepackage{booktabs}
\usepackage{tikz}
\usetikzlibrary{arrows.meta,positioning,calc}
\usepackage{hyperref}

% --- Paleta sobria ---
\definecolor{ink}{HTML}{18211D}
\definecolor{forest}{HTML}{123D31}
\definecolor{gold}{HTML}{A8843F}
\definecolor{warm}{HTML}{F7F5EF}
\definecolor{muted}{HTML}{6D746F}
\definecolor{soft}{HTML}{E8EAE4}
\definecolor{rust}{HTML}{8C493A}

\setbeamercolor{normal text}{fg=ink,bg=warm}
\setbeamercolor{structure}{fg=forest}
\setbeamercolor{frametitle}{fg=ink,bg=warm}
\setbeamercolor{title}{fg=ink,bg=warm}
\setbeamercolor{subtitle}{fg=muted,bg=warm}
\setbeamercolor{footline}{fg=muted,bg=warm}
\setbeamercolor{alerted text}{fg=rust}

\setbeamertemplate{navigation symbols}{}
\setbeamertemplate{itemize item}{\raise0.18ex\hbox{\scriptsize$\bullet$}}
\setbeamertemplate{itemize subitem}{\raise0.18ex\hbox{\tiny$\circ$}}
\setbeamertemplate{enumerate item}{\color{forest}\insertenumlabel.}
\setbeamersize{text margin left=0.78cm,text margin right=0.78cm}

% --- Tipografía ---
\setbeamerfont{title}{size=\fontsize{28}{31}\selectfont,series=\bfseries}
\setbeamerfont{subtitle}{size=\fontsize{14}{17}\selectfont}
\setbeamerfont{frametitle}{size=\fontsize{20}{23}\selectfont,series=\bfseries}
\setbeamerfont{normal text}{size=\fontsize{11}{14}\selectfont}

% --- Frametitle: aire + línea fina ---
\setbeamertemplate{frametitle}{%
  \vspace{0.28cm}
  \nointerlineskip
  \begin{beamercolorbox}[wd=\paperwidth,leftskip=0.78cm,rightskip=0.78cm,sep=0pt]{frametitle}
    \usebeamerfont{frametitle}\insertframetitle\par
    \vspace{0.10cm}
    {\color{gold}\rule{1.65cm}{0.65pt}}\par
  \end{beamercolorbox}
  \vspace{0.12cm}
}

% --- Pie ---
\setbeamertemplate{footline}{%
  \leavevmode%
  \hbox{%
    \begin{beamercolorbox}[wd=.82\paperwidth,ht=2.8ex,dp=1.2ex,leftskip=.78cm]{footline}%
      \scriptsize XXVII Semana de la Facultad de Economía · UAEMéx
    \end{beamercolorbox}%
    \begin{beamercolorbox}[wd=.18\paperwidth,ht=2.8ex,dp=1.2ex,rightskip=.78cm plus1fil]{footline}%
      \scriptsize\hfill \insertframenumber
    \end{beamercolorbox}%
  }%
}

% --- Macros de composición ---
\newcommand{\kicker}[1]{{\color{gold}\scriptsize\bfseries\MakeUppercase{#1}}}
\newcommand{\muted}[1]{{\color{muted}#1}}
\newcommand{\source}[1]{\vfill{\tiny\color{muted}Fuente: #1}}
\newcommand{\bignum}[2]{%
  {\fontsize{31}{33}\selectfont\bfseries\color{forest}#1}\par
  \vspace{0.05cm}{\small\color{muted}#2}}
\newcommand{\ruleline}{\par\vspace{0.18cm}\noindent{\color{soft}\rule{\linewidth}{0.55pt}}\par\vspace{0.18cm}}
\newcommand{\prompt}[1]{\texttt{\small #1}}

\hypersetup{
  colorlinks=true,
  linkcolor=forest,
  urlcolor=forest
}

% --- Datos ---
\title{La Inteligencia Artificial\\como aliada del emprendedor}
\subtitle{Herramientas prácticas para la nueva economía}
\author{Julio César Galindo}
\institute{XXVII Semana de la Facultad de Economía · Universidad Autónoma del Estado de México}
\date{26--30 de octubre de 2026}

\begin{document}

% ============================================================
% 1. PORTADA
% ============================================================
\begin{frame}[plain]
\vspace{0.45cm}
\kicker{Conferencia}
\vspace{0.65cm}

{\usebeamerfont{title}\color{ink}\inserttitle\par}
\vspace{0.36cm}
{\usebeamerfont{subtitle}\color{muted}\insertsubtitle\par}

\vfill
\begin{columns}[T,onlytextwidth]
  \begin{column}{0.68\textwidth}
    {\small\bfseries \insertauthor}\par
    \vspace{0.10cm}
    {\small\color{muted} \insertinstitute}
  \end{column}
  \begin{column}{0.32\textwidth}
    \raggedleft
    {\small\color{muted}\insertdate}\par
    \vspace{0.18cm}
    {\color{gold}\rule{2.55cm}{1.0pt}}
  \end{column}
\end{columns}
\vspace{0.25cm}
\end{frame}

% ============================================================
% 2. APERTURA INTERACTIVA
% ============================================================
\begin{frame}{Antes de hablar de IA, hablemos de tiempo}
\kicker{Interacción 01 · 60 segundos}
\vspace{0.35cm}

{\Large\bfseries ¿Qué tarea te roba más tiempo en tu negocio?}\par
\vspace{0.55cm}

\begin{columns}[T,onlytextwidth]
\begin{column}{0.23\textwidth}
\textbf{Clientes}\ruleline
\small Responder mensajes, cotizaciones y seguimiento.
\end{column}
\begin{column}{0.23\textwidth}
\textbf{Contenido}\ruleline
\small Redes, campañas, fotografías, textos y anuncios.
\end{column}
\begin{column}{0.23\textwidth}
\textbf{Administración}\ruleline
\small Facturas, gastos, reportes, minutas y archivos.
\end{column}
\begin{column}{0.23\textwidth}
\textbf{Decisiones}\ruleline
\small Precios, inventario, prioridades y planeación.
\end{column}
\end{columns}

\vfill
{\color{forest}\bfseries Levanta la mano por una sola opción.}\par\vspace{0.1cm}
\muted{La respuesta será nuestro mapa para la demostración.}
\end{frame}

% ============================================================
% 3. TESIS
% ============================================================
\begin{frame}{La tesis de esta charla}
\vspace{0.20cm}
\begin{center}
{\fontsize{25}{30}\selectfont\bfseries
La IA no sustituye un negocio.\\[0.12cm]
\color{forest}Reduce fricción alrededor del negocio.}
\end{center}

\vspace{0.65cm}
\begin{columns}[T,onlytextwidth]
\begin{column}{0.31\textwidth}
\kicker{Entrada}\par\vspace{0.18cm}
\textbf{Información}\par
\small\muted{mensajes, hojas de cálculo, documentos, imágenes, datos}
\end{column}
\begin{column}{0.06\textwidth}
\centering\vspace{0.45cm}{\Large\color{gold}$\longrightarrow$}
\end{column}
\begin{column}{0.26\textwidth}
\kicker{IA}\par\vspace{0.18cm}
\textbf{Transformación}\par
\small\muted{resumir, clasificar, redactar, comparar, proponer}
\end{column}
\begin{column}{0.06\textwidth}
\centering\vspace{0.45cm}{\Large\color{gold}$\longrightarrow$}
\end{column}
\begin{column}{0.31\textwidth}
\kicker{Salida}\par\vspace{0.18cm}
\textbf{Acción humana}\par
\small\muted{decidir, verificar, autorizar, negociar, crear confianza}
\end{column}
\end{columns}
\end{frame}

% ============================================================
% 4. QUÉ ES / NO ES
% ============================================================
\begin{frame}{IA generativa: una definición útil para negocios}
\begin{columns}[T,onlytextwidth]
\begin{column}{0.48\textwidth}
\kicker{Sí es}\par\vspace{0.25cm}
{\large\bfseries Un sistema que produce una respuesta a partir de contexto e instrucciones.}\par
\vspace{0.35cm}
\small Puede redactar, sintetizar, estructurar, clasificar, explicar y proponer alternativas.
\end{column}
\begin{column}{0.04\textwidth}
\centering{\color{soft}\rule{0.6pt}{4.7cm}}
\end{column}
\begin{column}{0.48\textwidth}
\kicker{No es}\par\vspace{0.25cm}
{\large\bfseries Una fuente infalible de verdad ni un responsable de tus decisiones.}\par
\vspace{0.35cm}
\small Puede equivocarse, completar huecos de forma plausible o perder contexto relevante.
\end{column}
\end{columns}

\vfill
{\color{forest}\bfseries Regla de oro:}\quad
\muted{delegar la primera versión; conservar el juicio final.}
\end{frame}

% ============================================================
% 5. MÉXICO
% ============================================================
\begin{frame}{La escala importa}
\kicker{México · Censos Económicos 2024}
\vspace{0.35cm}

\begin{columns}[T,onlytextwidth]
\begin{column}{0.38\textwidth}
\bignum{7.09 millones}{establecimientos en el país en 2024}
\end{column}
\begin{column}{0.38\textwidth}
\bignum{36.59 millones}{personas ocupadas en esos establecimientos}
\end{column}
\begin{column}{0.24\textwidth}
\vspace{0.10cm}
{\large\bfseries La pregunta no es solo}\par
\vspace{0.16cm}
{\large\color{forest}\bfseries “¿qué puede hacer la IA?”}\par
\vspace{0.16cm}
{\large\bfseries sino “¿dónde ahorra minutos cada día?”}
\end{column}
\end{columns}

\vfill
\small\muted{Los CE 2024 describen establecimientos y personal ocupado; la información económica corresponde principalmente a 2023.}
\source{INEGI, Resultados definitivos de los Censos Económicos 2024 (publicados en 2025).}
\end{frame}

% ============================================================
% 6. MAPA DE VALOR
% ============================================================
\begin{frame}{Cuatro lugares donde la IA puede crear valor}
\vspace{0.10cm}
\begin{center}
\begin{tikzpicture}[x=1cm,y=1cm]
  \draw[soft,line width=.7pt] (-5.6,0) -- (5.6,0);
  \draw[soft,line width=.7pt] (0,-2.2) -- (0,2.2);
  \node[align=left,text width=4.8cm] at (-3.05,1.15) {\kicker{01 · Ingresos}\\[-1mm]\Large\bfseries Vender mejor\\[1mm]\small\color{muted}propuestas, mensajes, ofertas, seguimiento};
  \node[align=left,text width=4.8cm] at (3.05,1.15) {\kicker{02 · Costos}\\[-1mm]\Large\bfseries Trabajar más rápido\\[1mm]\small\color{muted}documentos, reportes, clasificación, minutas};
  \node[align=left,text width=4.8cm] at (-3.05,-1.15) {\kicker{03 · Decisión}\\[-1mm]\Large\bfseries Ver patrones\\[1mm]\small\color{muted}comparaciones, escenarios, anomalías, prioridades};
  \node[align=left,text width=4.8cm] at (3.05,-1.15) {\kicker{04 · Experiencia}\\[-1mm]\Large\bfseries Atender mejor\\[1mm]\small\color{muted}respuestas consistentes, FAQ, personalización};
\end{tikzpicture}
\end{center}
\end{frame}

% ============================================================
% 7. FINANZAS
% ============================================================
\begin{frame}{Caso 1 · Finanzas: de datos a decisiones}
\begin{columns}[T,onlytextwidth]
\begin{column}{0.47\textwidth}
\kicker{Flujo de trabajo}\par\vspace{0.25cm}
\begin{enumerate}
\item Exportar movimientos de banco o ventas.
\item Quitar datos sensibles innecesarios.
\item Pedir clasificación y resumen semanal.
\item Revisar excepciones y confirmar cifras.
\item Convertir el resultado en una decisión.
\end{enumerate}
\end{column}
\begin{column}{0.06\textwidth}
\centering{\color{soft}\rule{0.6pt}{5.2cm}}
\end{column}
\begin{column}{0.47\textwidth}
\kicker{Ejemplo de instrucción}\par\vspace{0.25cm}
\prompt{Analiza esta tabla de ingresos y gastos.}\par
\prompt{1) clasifica por categoría;}\par
\prompt{2) señala variaciones >20\%;}\par
\prompt{3) detecta cargos atípicos;}\par
\prompt{4) no inventes datos faltantes;}\par
\prompt{5) termina con 3 preguntas de control.}
\vspace{0.40cm}

{\small\color{rust}\bfseries Nunca confundas una explicación convincente con una conciliación contable.}
\end{column}
\end{columns}
\end{frame}

% ============================================================
% 8. PROMPT INTERACTIVO
% ============================================================
\begin{frame}{El prompt importa menos que el contexto bien puesto}
\kicker{Interacción 02 · Mejoremos una instrucción}
\vspace{0.30cm}

\only<1>{
{\large\bfseries Prompt débil}\par\vspace{0.22cm}
\prompt{“Hazme una campaña para vender más.”}
\vspace{0.65cm}

{\large\bfseries ¿Qué información falta?}\par
\small\muted{Pide al público tres datos antes de avanzar.}
}

\only<2>{
{\large\bfseries Una versión mucho más útil}\par\vspace{0.25cm}
\prompt{“Tengo una cafetería independiente en Toluca.}\par
\prompt{Mi ticket promedio es de \$115 y quiero aumentar visitas}\par
\prompt{de lunes a jueves entre 16:00 y 19:00. Mi público}\par
\prompt{principal son estudiantes y profesionistas de 20–35 años.}\par
\prompt{Propón 3 campañas de bajo costo, cada una con oferta,}\par
\prompt{mensaje, canal, métrica y riesgo. No uses descuentos >15\%.”}
\vspace{0.45cm}

{\color{forest}\bfseries Mejor contexto $\rightarrow$ mejores alternativas $\rightarrow$ mejor juicio humano.}
}
\end{frame}

% ============================================================
% 9. MARKETING
% ============================================================
\begin{frame}{Caso 2 · Marketing: del “publica algo” al experimento}
\begin{columns}[T,onlytextwidth]
\begin{column}{0.32\textwidth}
\kicker{Hipótesis}\par\vspace{0.15cm}
\textbf{Una audiencia}\par
\small\muted{¿Para quién es? ¿Qué problema intenta resolver?}
\end{column}
\begin{column}{0.32\textwidth}
\kicker{Mensaje}\par\vspace{0.15cm}
\textbf{Una promesa verificable}\par
\small\muted{¿Qué cambia para el cliente? ¿Qué evidencia la sostiene?}
\end{column}
\begin{column}{0.32\textwidth}
\kicker{Medición}\par\vspace{0.15cm}
\textbf{Una métrica}\par
\small\muted{clics, citas, ventas, recompra, tiempo de respuesta}
\end{column}
\end{columns}

\vspace{0.65cm}
{\Large\bfseries La IA puede multiplicar variantes.}\par
\vspace{0.10cm}
{\Large\color{forest}\bfseries El emprendedor decide cuál merece una prueba real.}
\end{frame}

% ============================================================
% 10. OPERACIONES
% ============================================================
\begin{frame}{Caso 3 · Operaciones: de experiencia a proceso}
\begin{columns}[T,onlytextwidth]
\begin{column}{0.41\textwidth}
\kicker{Antes}\par\vspace{0.20cm}
{\large\bfseries “Solo Ana sabe cómo se hace.”}\par
\vspace{0.24cm}
\small La operación depende de memoria, mensajes dispersos y correcciones repetidas.
\end{column}
\begin{column}{0.18\textwidth}
\centering\vspace{0.65cm}
{\fontsize{25}{25}\selectfont\color{gold}$\longrightarrow$}\par
\small\muted{documentar}
\end{column}
\begin{column}{0.41\textwidth}
\kicker{Después}\par\vspace{0.20cm}
{\large\bfseries “Existe un procedimiento revisable.”}\par
\vspace{0.24cm}
\small La IA ayuda a ordenar notas, crear checklist, detectar ambigüedades y preparar capacitación.
\end{column}
\end{columns}

\vfill
{\color{forest}\bfseries Uso ideal:}\quad \muted{primera versión del procedimiento + revisión de quien realmente hace el trabajo.}
\end{frame}

% ============================================================
% 11. AUTOMATIZACIÓN
% ============================================================
\begin{frame}{Automatizar no significa quitar al humano}
\vspace{0.15cm}
\begin{columns}[T,onlytextwidth]
\begin{column}{0.22\textwidth}
\centering\kicker{Disparador}\par\vspace{0.10cm}\bfseries llega un correo
\end{column}
\begin{column}{0.04\textwidth}
\centering\vspace{0.18cm}{\Large\color{gold}$\rightarrow$}
\end{column}
\begin{column}{0.22\textwidth}
\centering\kicker{IA}\par\vspace{0.10cm}\bfseries clasifica y resume
\end{column}
\begin{column}{0.04\textwidth}
\centering\vspace{0.18cm}{\Large\color{gold}$\rightarrow$}
\end{column}
\begin{column}{0.22\textwidth}
\centering\kicker{Regla}\par\vspace{0.10cm}\bfseries decide la ruta
\end{column}
\begin{column}{0.04\textwidth}
\centering\vspace{0.18cm}{\Large\color{gold}$\rightarrow$}
\end{column}
\begin{column}{0.22\textwidth}
\centering\kicker{Acción}\par\vspace{0.10cm}\bfseries crea borrador
\end{column}
\end{columns}

\vspace{0.75cm}
\begin{columns}[T,onlytextwidth]
\begin{column}{0.48\textwidth}
{\large\bfseries Automático}\par
\small\muted{clasificar, extraer, resumir, preparar, registrar}
\end{column}
\begin{column}{0.48\textwidth}
{\large\bfseries Bajo autorización}\par
\small\muted{enviar, pagar, publicar, comprometer precio, tratar datos sensibles}
\end{column}
\end{columns}

\vfill
{\color{rust}\bfseries El punto de control humano debe estar antes de la acción irreversible.}
\end{frame}

% ============================================================
% 12. DEMO
% ============================================================
\begin{frame}[label=demo]{Demostración: del problema al prototipo}
\kicker{Interacción 03 · 8 minutos}
\vspace{0.25cm}

\begin{columns}[T,onlytextwidth]
\begin{column}{0.56\textwidth}
{\large\bfseries El público elige un negocio:}\par
\small cafetería · consultorio · tienda · despacho · taller
\vspace{0.32cm}

{\large\bfseries Y una tarea:}\par
\small responder clientes · cotizar · resumir ventas · crear campaña · documentar proceso
\vspace{0.42cm}

{\color{forest}\bfseries En vivo construimos:}\par
\small contexto $\rightarrow$ instrucción $\rightarrow$ primera salida $\rightarrow$ crítica $\rightarrow$ versión útil
\end{column}
\begin{column}{0.44\textwidth}
\kicker{Pregunta al público}\par\vspace{0.18cm}
{\large\bfseries “¿Qué dato cambiaría por completo la respuesta?”}\par
\vspace{0.38cm}
\small\muted{La interacción enseña algo central: usar IA es diseñar una conversación, no presionar un botón.}
\end{column}
\end{columns}

\vfill
{\scriptsize\href{https://chatgpt.com/}{Abrir ChatGPT para la demostración} \quad·\quad Llevar un ejemplo alternativo preparado por si falla la conexión.}
\vspace{0.35cm}
\end{frame}

% ============================================================
% 13. ARQUITECTURA DE PROMPT
% ============================================================
\begin{frame}{Una arquitectura mejor que los “prompts mágicos”}
\vspace{0.10cm}
\begin{enumerate}
\item \textbf{Contexto}\;— negocio, público, objetivo, datos disponibles.
\item \textbf{Tarea}\;— qué debe producir exactamente.
\item \textbf{Restricciones}\;— presupuesto, tono, límites, políticas, fechas.
\item \textbf{Formato}\;— tabla, correo, checklist, guion, JSON, pasos.
\item \textbf{Verificación}\;— qué debe marcar como incierto o revisar antes de concluir.
\end{enumerate}

\vspace{0.55cm}
{\small\prompt{“Cuando falte información, no la inventes: indícala como ‘dato no disponible’ y formula la pregunta necesaria.”}}

\vfill
{\color{forest}\bfseries No busques una frase perfecta. Diseña un proceso reproducible.}
\end{frame}

% ============================================================
% 14. DATOS Y PRIVACIDAD
% ============================================================
\begin{frame}{Antes de copiar y pegar: una pausa de 10 segundos}
\begin{columns}[T,onlytextwidth]
\begin{column}{0.49\textwidth}
\kicker{Evita introducir sin necesidad}\par\vspace{0.20cm}
\small
\begin{itemize}
\item contraseñas, llaves y credenciales;
\item datos bancarios completos;
\item expedientes o datos personales sensibles;
\item información confidencial de clientes o proveedores;
\item secretos comerciales no autorizados.
\end{itemize}
\end{column}
\begin{column}{0.02\textwidth}
\centering{\color{soft}\rule{0.6pt}{5.0cm}}
\end{column}
\begin{column}{0.49\textwidth}
\kicker{Mejor práctica}\par\vspace{0.20cm}
\small
\begin{itemize}
\item anonimiza nombres e identificadores;
\item comparte solo lo necesario;
\item revisa políticas y configuración del servicio;
\item usa cuentas y herramientas aprobadas por la organización;
\item conserva control de documentos originales.
\end{itemize}
\end{column}
\end{columns}

\vfill
{\color{rust}\bfseries Conveniencia no equivale a autorización.}
\end{frame}

% ============================================================
% 15. ERRORES
% ============================================================
\begin{frame}{Una respuesta convincente puede ser incorrecta}
\kicker{Control de calidad}
\vspace{0.25cm}

\begin{center}
\begin{tikzpicture}
\node[anchor=north west,align=left,text width=3.75cm] at (-5.55,1.15) {\textbf{1. Hechos}\\[2mm]\small Contrastar nombres, cifras, fechas y normas con una fuente primaria.};
\node[anchor=north west,align=left,text width=3.75cm] at (-1.80,1.15) {\textbf{2. Cálculos}\\[2mm]\small Reproducir la cuenta en Excel, Python o una calculadora.};
\node[anchor=north west,align=left,text width=3.75cm] at (1.95,1.15) {\textbf{3. Decisión}\\[2mm]\small Definir quién asume la responsabilidad si la recomendación falla.};
\draw[soft,line width=.6pt] (-5.55,0.72)--(-2.05,0.72);
\draw[soft,line width=.6pt] (-1.80,0.72)--(1.70,0.72);
\draw[soft,line width=.6pt] (1.95,0.72)--(5.45,0.72);
\end{tikzpicture}
\end{center}

\vspace{0.72cm}
{\Large\bfseries Flujo seguro:}\quad
{\Large\color{forest}\bfseries generar $\rightarrow$ comprobar $\rightarrow$ decidir.}
\end{frame}

% ============================================================
% 16. ROI
% ============================================================
\begin{frame}{¿Vale la pena? Midámoslo}
\vspace{0.15cm}
\begin{center}
{\fontsize{20}{24}\selectfont
$\displaystyle \text{Valor mensual} \approx
(\text{horas ahorradas})(\text{valor por hora})
+ \text{margen incremental}
- \text{costo de herramientas}
- \text{costo de revisión}$}
\end{center}

\vspace{0.65cm}
\begin{columns}[T,onlytextwidth]
\begin{column}{0.30\textwidth}
\kicker{Mide}\par\vspace{0.12cm}
\textbf{Tiempo}\par
\small\muted{minutos antes vs. después}
\end{column}
\begin{column}{0.30\textwidth}
\kicker{Mide}\par\vspace{0.12cm}
\textbf{Calidad}\par
\small\muted{errores, retrabajo, respuesta}
\end{column}
\begin{column}{0.30\textwidth}
\kicker{Mide}\par\vspace{0.12cm}
\textbf{Resultado}\par
\small\muted{ventas, citas, cobranza, recompra}
\end{column}
\end{columns}

\vfill
{\color{forest}\bfseries Una automatización que no se mide se convierte fácilmente en entretenimiento tecnológico.}
\end{frame}

% ============================================================
% 17. PLAN 30 DÍAS
% ============================================================
\begin{frame}{Un plan de 30 días: pequeño, medible y reversible}
\begin{columns}[T,onlytextwidth]
\begin{column}{0.23\textwidth}
\kicker{Semana 1}\par\vspace{0.15cm}
\textbf{Elegir}\par
\small una tarea repetitiva, frecuente y de bajo riesgo.
\end{column}
\begin{column}{0.23\textwidth}
\kicker{Semana 2}\par\vspace{0.15cm}
\textbf{Prototipar}\par
\small un flujo simple con ejemplos reales anonimizados.
\end{column}
\begin{column}{0.23\textwidth}
\kicker{Semana 3}\par\vspace{0.15cm}
\textbf{Medir}\par
\small tiempo, calidad, errores y efecto en clientes.
\end{column}
\begin{column}{0.23\textwidth}
\kicker{Semana 4}\par\vspace{0.15cm}
\textbf{Decidir}\par
\small mantener, ajustar, automatizar más o descartar.
\end{column}
\end{columns}

\vspace{0.80cm}
{\Large\bfseries Empieza donde el error sea barato}\par
\vspace{0.08cm}
{\Large\color{forest}\bfseries y donde aprender sea valioso.}
\end{frame}

% ============================================================
% 18. HUMANO
% ============================================================
\begin{frame}{Lo que se vuelve más valioso cuando la IA abarata la producción}
\vspace{0.20cm}
\begin{center}
\begin{tikzpicture}
\node[align=center] at (-4.6,1.1) {\kicker{Juicio}\\\Large\bfseries decidir};
\node[align=center] at (-1.55,1.1) {\kicker{Contexto}\\\Large\bfseries entender};
\node[align=center] at (1.55,1.1) {\kicker{Confianza}\\\Large\bfseries relacionar};
\node[align=center] at (4.6,1.1) {\kicker{Responsabilidad}\\\Large\bfseries responder};
\draw[gold,line width=.8pt] (-5.3,0.25) -- (5.3,0.25);
\node[align=center,text width=10.8cm] at (0,-1.05) {\Large\color{forest}\bfseries La ventaja no será “tener IA”.\\[-1mm]\Large\bfseries Será integrarla a un negocio que ya sabe qué problema quiere resolver.};
\end{tikzpicture}
\end{center}
\end{frame}

% ============================================================
% 19. CIERRE
% ============================================================
\begin{frame}{Una última pregunta}
\kicker{Interacción final}
\vspace{0.45cm}

\begin{center}
{\fontsize{24}{29}\selectfont\bfseries
¿Qué tarea de tu próximo lunes\\
\color{forest}no debería seguir haciéndose igual?}
\end{center}

\vspace{0.50cm}
\begin{center}
\small\muted{Escríbela en una frase. Esa frase puede convertirse en tu primer experimento de IA.}
\end{center}

\vspace{0.42cm}
{\color{gold}\rule{\linewidth}{0.8pt}}\par
\vspace{0.22cm}
{\large\bfseries Gracias.}\quad
\muted{Preguntas, objeciones y casos reales.}\hfill
{\small jcgalindo@up.edu.mx · ivlicaesar.github.io}
\end{frame}

\end{document}
